\uuid{H7AB}
\exo7id{1812}
\titre{exo7 1812}
\auteur{drutu}
\organisation{exo7}
\datecreate{2003-10-01}
\isIndication{false}
\isCorrection{false}
\chapitre{Fonction de plusieurs variables}
\sousChapitre{Différentiabilité}
\cpge{mpsi-mp2i.fonctions-deux-variables.a, pcsi.fonctions-deux-variables.a, ptsi.fonctions-deux-variables.a, mpsi-mp2i.fonctions-deux-variables.b, pcsi.fonctions-deux-variables.b, ptsi.fonctions-deux-variables.b, bcpst1.fonctions-deux-variables.b, bcpst1.fonctions-deux-variables.c, mp-mpi.calcul-differentiel-optimisation.b, pc.fonctions-plusieurs-variables.c, psi.fonctions-plusieurs-variables.c}
\module{Analyse}
\niveau{L2}
\difficulte{}

\contenu{
\question{
Soit $f:\R^2 \to \R $, 
$$
f(x,y)=\left\{ 
        \begin{array}{cc}
        \frac{x^2y+xy^2}{x^2+y^2} & \mbox{ si }(x,y)\neq (0,0) \\
        0 & \mbox{ si }(x,y)=(0,0) 
        \end{array}
        \right . 
$$
Montrer que $f$ est continue en $(0,0)$ et admet des d\'eriv\'ees partielles 
dans toutes les directions, mais n'y est pas diff\'erentiable.
}
}
